\section{The inherited selected short--short interface}
\label{sec:interface}

Fix \(h\in\Z\setminus\{0\}\), put \(H=\rad|h|\), and let
\begin{equation}
 R=\lfloor X^{1/2-\delta}\rfloor,\qquad 0<\delta<\frac12,
 \label{eq:R-def}
\end{equation}
with \(R\ge H^2\).  In one separated row packet,
\begin{equation}
 Q=LD,\qquad QJ\asymp X,\qquad L>\max\{2R,2|h|\},
 \label{eq:row-scales}
\end{equation}
and
\begin{equation}
 \omega=(\ell,d),\quad \ell\in[L/2,2L]\text{ prime},\quad
 d\in[D,2D],\quad d\le R^{1/2},\quad m_\omega=\ell d\asymp Q.
 \label{eq:rows}
\end{equation}
The large prime is unique, so \(m_\omega\) determines the row label.
The bounded separated weights satisfy
\begin{equation}
 \|x\|_\infty+\|y\|_\infty\ll_\eps X^\eps,\qquad
 \sum_m|x_m|^2+\sum_n|y_n|^2\ll QX^\eps.
 \label{eq:row-energy}
\end{equation}

The actual generic mask is
\begin{equation}
 \mathfrak m(m,n)=\one_{\ell_m\ne\ell_n}
 \one_{|m-n|>QX^{-\kappa_0}}
 \one_{(d_m,d_n)\le X^{\kappa_0}},
 \label{eq:generic-mask}
\end{equation}
where \(\kappa_0>0\) is fixed.  We use only
\begin{equation}
 |\mathfrak m(m,n)|\le1,\qquad \mathfrak m(m,m)=0,\qquad
 \mathfrak m\text{ is independent of }u,v,q.
 \label{eq:mask-properties}
\end{equation}
Thus the theorem applies to any fixed bounded off-diagonal row-pair mask
with these properties, but not as stated to a modulus-dependent mask.

The divisor coefficient is either
\begin{equation}
 \vartheta_R(u)=\lambda'_R(u)
 \quad\text{or}\quad \vartheta_R(u)=b_R(u)\one_{u\le R}.
 \label{eq:theta-choices}
\end{equation}
Both are real, squarefree-supported on \(u\le R\), and
\begin{equation}
 |\vartheta_R(u)|\ll_\eps X^\eps.
 \label{eq:theta-bound}
\end{equation}

For squarefree \(u,v\le R\), put \(g=(u,v)\), \(q=[u,v]\), and retain
only \((q,H)=1\).  A row pair is admissible when
\begin{equation}
 (u,mH)=(v,nH)=1,\qquad m\equiv n\pmod g.
 \label{eq:pair-admissibility}
\end{equation}
It determines a unique \(\xi\in\cU_q=(\Z/q\Z)^\times\) through
\begin{equation}
 m\xi\equiv-h\pmod u,\qquad n\xi\equiv-h\pmod v.
 \label{eq:joint-class}
\end{equation}
Define
\begin{align}
 Z_{u,v}(\xi)&=\sum_{m,n}x_my_n\mathfrak m(m,n)
 \one_{m\equiv-h\overline\xi\ (u)}
 \one_{n\equiv-h\overline\xi\ (v)},
 \label{eq:pair-fiber}\\
 Z_q(\xi)&=\sum_{\substack{u,v\le R\\\operatorname{lcm}(u,v)=q}}
 \vartheta_R(u)\vartheta_R(v)Z_{u,v}(\xi).
 \label{eq:lcm-fiber}
\end{align}
Inadmissible row--modulus combinations contribute zero.

Let
\begin{equation}
 \overline Z_{u,v}=\frac1{\varphi(q)}\sum_{\xi\in\cU_q}Z_{u,v}(\xi),
 \qquad D_{u,v}=Z_{u,v}-\overline Z_{u,v},
 \label{eq:pair-centering}
\end{equation}
and define \(\overline Z_q,D_q\) similarly.  Centering is linear:
\begin{equation}
 D_q=\sum_{[u,v]=q}\vartheta_R(u)\vartheta_R(v)D_{u,v}.
 \label{eq:centered-pair-sum}
\end{equation}
This licenses pairwise triangle inequalities, not orthogonality.

We use the normalizations
\begin{equation}
 e(t)=e^{2\pi i t},\qquad e_q(t)=e(t/q),\qquad
 c_s(r)=\sum_{\xi\in(\Z/s\Z)^\times}e_s(r\xi),
 \label{eq:additive-normalizations}
\end{equation}
and, for a function \(W\) on \(\cU_q\),
\begin{equation}
 \|W\|_{2,\cU_q}^2
 :=\sum_{\xi\in\cU_q}|W(\xi)|^2.
 \label{eq:unnormalized-fiber-norm}
\end{equation}

For \(F\in C_c^\infty(\R)\), with
\(\widehat F(t)=\int_\R F(x)e(-tx)\,dx\), put
\begin{equation}
 \cL_q=\frac{J}{Hq}\sum_{r\ne0}c_H(r)
 \widehat F\left(\frac{rJ}{Hq}\right)
 \sum_{\xi\in\cU_q}Z_q(\xi)e_q(r\overline H\xi).
 \label{eq:scalar-packet}
\end{equation}
We write \(\cL_q^{\rm mean}\) and \(\cL_q^{\rm disc}\) for the same
linear functional with \(Z_q\) replaced by \(\overline Z_q\) and
\(D_q\), respectively; hence
\(\cL_q=\cL_q^{\rm mean}+\cL_q^{\rm disc}\).
TPC-21 proves for its mean part
\begin{equation}
 \sum_q|\cL_q^{\rm mean}|\ll_\eps Q^2X^\eps.
 \label{eq:imported-mean-bound}
\end{equation}
For the discrepancy, TPC-20 supplies under its no-wrap and polynomial
tail-mass hypotheses
\begin{equation}
 \sum_{q\in\cQ}|\cL_q^{\rm disc}|
 \ll_{H,F,\eps}X^{\eps/2}\sqrt{\frac JH}
 \sum_{q\in\cQ}\lVert D_q\rVert_{2,\cU_q}
 +O_A(X^{-A}\cZ_1).
 \label{eq:inherited-norm-route}
\end{equation}
Here \(\cZ_1\) is the total absolute coefficient mass of the separated
packet.  The inherited polynomial tail-mass hypothesis includes
\(\cZ_1=X^{O(1)}\), so the last term is smaller than any prescribed
fixed power after choosing \(A\).
The natural scale is
\begin{equation}
 Q^2J\asymp XQ.
 \label{eq:natural-scale}
\end{equation}

We absorb dyadic decompositions and
\begin{equation}
 \tau(n)^C+3^{\omega(n)}\ll_{C,\eps}X^\eps
 \qquad(n\le X^{O(1)})
 \label{eq:divisor-soft}
\end{equation}
into \(X^\eps\).  A conductor block of size one contains exact
conductor one.

\begin{theorem}[Selected short--short connected closure]
\label{thm:main-closure}
Assume explicitly, for some fixed \(\eps_0>0\), the TPC-20 no-wrap
condition \(J\ge4HX^{\eps_0}\) and its polynomial fiber-tail mass hypothesis
used in \eqref{eq:inherited-norm-route}.  Write
\(Q=X^{\beta+o(1)}\), \(J=X^{1-\beta+o(1)}\).  Suppose either
\begin{equation}
 \beta<2\delta
 \label{eq:occupancy-region}
\end{equation}
with a fixed margin, or
\begin{equation}
 1-4\delta<\beta<\min\left\{\frac23,6\delta\right\}
 \label{eq:main-hybrid-region}
\end{equation}
with fixed margins.  Assume also the strict TPC-20 single-leg condition
\begin{equation}
 \beta<\frac12+\delta
 \label{eq:main-single-leg}
\end{equation}
with a fixed margin.  Then some \(\eta>0\) satisfies, uniformly for
every dyadic \(Y\le R^2\),
\begin{equation}
 \sum_{\substack{q\asymp Y\\q\text{ squarefree}\\(q,H)=1}}
 |\cL_q^{\rm disc}|\ll_{F,H,\eps}XQ\,X^{-\eta+\eps}.
 \label{eq:main-discrepancy-bound}
\end{equation}
After the inherited dyadic and separated-packet reassembly, this and
\eqref{eq:imported-mean-bound} close the selected short--short packet.
\end{theorem}

At \(\beta=2\delta\), occupancy has no fixed-power margin on
\(Y\asymp R^2\).  The second alternative is the new result.  Neither
alternative includes the omitted large-divisor residual legs.
